\section{The model}\label{sec:prelim}

Let $F$ be a field and $d\ge1$. Rows and columns of a matrix $X\in\Mat_d(F)$ are indexed by $\Z/d$, and $x_c\in F^{\Z/d}$
is row~$c$. The rotation $\sigma^c$ moves the columns: $(\sigma^cX)_{i,j}=X_{i,j+c}$ (\lean{SwapRigidity.rot}). We use
the parent paper's definitions.

\begin{definition}[Class-uniform rules {\cite[Definition~6.3]{SeventeenthPaper2026}}]\label{def:class-uniform}
\formal{ClassUniform.lean}{SwapRigidity.rot,SwapRigidity.IsClassUniformCode}
Fix a colouring of $\Z/d$, that is, a partition $\mathcal P$ into classes. A \emph{class-uniform rule} is a family of
linear maps $\cA_J\colon\Mat_d(F)\to F^{\Z/d}$, one for each class $J$. A subspace $\X\subseteq\Mat_d(F)$ is a
\emph{code} of the rule if $x_c=\cA_J(\sigma^cX)$ for every $X\in\X$, every class $J$ and every $c\in J$. We write
$D(\mathcal P)$ for the largest dimension of a code of a class-uniform rule for $\mathcal P$, over the field at hand.
\end{definition}

The \emph{code of} a rule, as opposed to a code of it, is the space of all matrices that satisfy its equations, its
largest code. The file named in the status line is the parent paper's \leanfile{ClassUniform.lean}. In Lean the colouring is a map
\texttt{cls} from $\Z/d$ to a type of classes, and \lean{SwapRigidity.IsClassUniformCode} asserts that decoders exist.

\begin{theorem}[Theorem~TP {\cite[Theorem~6.4]{SeventeenthPaper2026}}]\label{thm:TP}
\cited{the parent paper, Theorem~6.4, Lean-verified there (\lean{SwapRigidity.IsClassUniformCode.finrank_le},
\lean{SwapRigidity.IsClassUniformCode.two_mul_finrank_le}).}
Let $\X$ be a code of a class-uniform rule.
\begin{enumerate}[label=(\alph*),leftmargin=2em]
\item Let $\Delta_0\subseteq\Z/d$ be a set of $s+1\ge2$ shifts, and let $\chi$ be the number of components of the graph
on $\Z/d$ in which $c\sim c'$ when $c$ and $c'$ lie in one class and $c'-c\in\Delta_0-\Delta_0$. Then
$(s+1)\dim\X\le s\,d^2+d\chi$.
\item Let $S\subseteq\Z/d$ be a non-empty set that contains every $c$ whose class differs from that of $c-1$. Then
$2\dim\X\le d^2+d\card S$. In particular, if the classes are $M\ge2$ cyclic intervals, then $\dim\X\le(d^2+Md)/2$.
\end{enumerate}
\end{theorem}

For $d\ge2$, part~(b) follows from~(a) with $\Delta_0=\{0,1\}$, whose components are the maximal runs of one class.
In Lean, (a) takes in place of $\chi$ the size of any set of rows that meets every component, and (b) a non-empty set of
rows that contains every class change. For the parent paper a
partition into $M$ classes has an \emph{almost-period} if some pair of shifts $\Delta_0=\{0,\delta\}$ leaves only $M$
components in the graph of~(a); intervals have the almost-period $1$. The classes $\{0,1,3\}$ and
$\{2,4,5\}$ of $\Z/6$ have none: every nonzero shift leaves four components
\cite[Remark~8.1(c)]{SeventeenthPaper2026}.

\begin{proposition}[{\cite[Proposition~8.2]{SeventeenthPaper2026}}]\label{prop:F2}
\cited{the parent paper, Proposition~8.2, Lean-verified there (\lean{SwapRigidity.isClassUniformCode_aperiodicCode},
\lean{SwapRigidity.finrank_aperiodicCode}, \lean{SwapRigidity.aperiodicCls_eq_true_iff}).}
Over $\F_2$ the classes $\{0,1,3\}$ and $\{2,4,5\}$ of $\Z/6$ carry a code of a class-uniform rule of dimension $25$,
while $(d^2+Md)/2=24$.
\end{proposition}

The parent paper asks~\cite[Problem~9.1]{SeventeenthPaper2026}: (a) whether $D(\mathcal P)\le(\frac12+o(1))d^2$ for two
classes, or for $o(d)$ classes; (b) the value of $D$ for these two classes over $\F_2$, which is between $25$ and $28$;
and (c) for a family of partitions and codes, with a description that can be analysed, in which the \emph{excess}
$D(\mathcal P)-(d^2+Md)/2$ grows; its later revision~\cite{SeventeenthPaperRevised2026} states~(c) for a fixed
number of classes, as \cref{prob:bounded} does. We call $\dim\X-(d^2+Md)/2$ the excess of a code $\X$; the excess of a partition is at
least that of each of its codes.
